\documentclass[11pt,a4paper]{article}
\usepackage[T2A]{fontenc}
\usepackage[utf8]{inputenc}
\usepackage[russian,english]{babel}
\usepackage{amsmath,amssymb,amsthm,mathtools}
\usepackage{setspace}
\usepackage{enumitem}
\usepackage[hidelinks]{hyperref}
\setstretch{1.3}
\frenchspacing
\hyphenpenalty=10000
\exhyphenpenalty=10000
\renewcommand{\leq}{\leqslant}
\renewcommand{\le}{\leqslant}
\renewcommand{\geq}{\geqslant}
\renewcommand{\ge}{\geqslant}
\newcommand{\head}[2][0pt]{\par\vspace{\dimexpr0.85em+#1\relax}\noindent\textbf{\underline{#2}}\quad}
\begin{document}
\begin{center}
{\huge\bfseries Limits -- 3\par}\vspace{1.07em}
{\Large Asymptotics of sequences\par}\vspace{1.1em}
{\foreignlanguage{russian}{1 сентября 2026 года}\par}
\end{center}\vspace{0pt}

\noindent\textbf{\underline{Theoretical minimum.}}

\vspace{5pt}\noindent\textbf{\underline{Def.}}\quad Knowing the limit is not enough: an olympiad
problem usually asks for the rate. The standard answer has the form $x_n\sim Cn^{-\alpha}$ or
$x_n=L+Cq^n+o(q^n)$, and the choice between the two is decided by the behaviour of $f$ at the
fixed point.

\head{Theorem 1 (neutral fixed point).} Let $x_n\to0$, $x_n>0$ and
$x_{n+1}=x_n-ax_n^p+o(x_n^p)$ with $a>0$, $p>1$. Then
\[
x_n\sim\Bigl(\frac1{a(p-1)n}\Bigr)^{1/(p-1)}.
\]
In particular $x_{n+1}=x_n-x_n^2$ gives $nx_n\to1$, while $x_{n+1}=\sin x_n$ (where $p=3$,
$a=\frac16$) gives $x_n\sim\sqrt{3/n}$.

\head[5.8pt]{Theorem 2 (Stolz--Ces\`aro).} If $b_n$ is strictly increasing and unbounded and
$\lim\frac{a_{n+1}-a_n}{b_{n+1}-b_n}=L$, then $\lim\frac{a_n}{b_n}=L$. This is the discrete
analogue of l'H\^opital's rule, and it is exactly what proves Theorem 1: apply it to
$y_n=x_n^{1-p}$.

\head[3.5pt]{Fact 1.} The telescoping trick in powers of $x_n$: from $x_{n+1}=x_n-ax_n^p$ one gets
$x_{n+1}^{1-p}-x_n^{1-p}\to a(p-1)$, and then Stolz finishes the job. This trick solves almost
every problem of the sheet.

\noindent\rule{\textwidth}{0.4pt}\par\vspace{8.4pt}
\begin{enumerate}
\item (IMC 2014) Consider the sequence
$(a_n)_{n\ge1}=(1,1,2,1,2,3,1,2,3,4,1,2,3,4,5,1,\dots)$.
Find all pairs $(\alpha,\beta)$ of positive reals such that
$\lim_{n\to\infty}\frac{\sum_{k=1}^na_k}{n^\alpha}=\beta$.

\item (Putnam 1966) Let $0<x_0<1$ and $x_{n+1}=x_n-x_n^2$ for $n\ge0$. Compute
$\lim_{n\to\infty}nx_n$.

\newpage
\item For an arbitrary $x_0\in(0,\pi)$ define $x_{n+1}=\sin x_n$. Compute
$\lim_{n\to\infty}\sqrt n\,x_n$, and find the next term of the asymptotics.

\item (AKHIMO 2022) Let $x_1\in\bigl(0,\frac\pi2\bigr)$ and
$x_{n+1}=\sqrt[3]{\frac{\sin^3x_1+\dots+\sin^3x_n}n}$ for $n\ge1$. Find
$\lim_{n\to\infty}x_n^2\ln n$.

\item Let $f\colon[a,b]\to[a,b]$ be an increasing function. Show that there exists $\xi\in[a,b]$
with $f(\xi)=\xi$, and give an example of a decreasing $f$ for which the iterations
$x_{n+1}=f(x_n)$ do not converge although the fixed point is unique.
\end{enumerate}
\end{document}
